\documentclass[11pt,a4paper]{article}
\usepackage[utf8]{inputenc}
\usepackage[spanish]{babel}
\usepackage{amsmath}
\usepackage{amsfonts}
\usepackage{amssymb}
\usepackage{geometry}
\usepackage{hyperref}

\geometry{left=2.5cm, right=2.5cm, top=3cm, bottom=3cm}

\title{\textbf{Desarrollo Matemático y Justificación del Cambio de Variable $Q = h \phi_s$\\ en Dunas Bidispersas}}
\author{Felipe Espinoza}
\date{\today}

\begin{document}

\maketitle

\section{Introducción}
En la modelación de dunas bidispersas que migran sobre un lecho, la topografía del fondo $h(x,t)$ varía en el espacio y en el tiempo. Para resolver las ecuaciones de advección-segregación-difusión de manera eficiente, se realiza una transformación a coordenadas adaptadas al contorno (coordenadas $\eta$ o tipo sigma). 

En este documento se presenta de manera detallada la justificación física y la derivación matemática del cambio de variable:
\begin{equation}
    Q = h(x,t) \, \phi_s
\end{equation}
donde $\phi_s$ es la concentración volumétrica (fracción adimensional) de partículas finas y $Q$ es el espesor equivalente de finos (medido en metros).

\section{Justificación Física: Conservación Global de la Masa}
Considere un dominio bidimensional $(x,z)$ donde la base del lecho está fija en $z = H_{base} = 0$ y la superficie de la duna está dada por $z = h(x,t)$. La masa total (o volumen neto) de partículas finas $M_{\text{finos}}$ en el sistema está dada por la integral:
\begin{equation}
    M_{\text{finos}} = \int_{0}^{L} \int_{0}^{h(x,t)} \phi_s(x, z, t) \, dz \, dx
\end{equation}

Para resolver el problema sobre una malla rectangular computacional fija, realizamos una transformación de la coordenada vertical:
\begin{equation}
    \eta = \frac{z}{h(x,t)} \implies z = \eta \cdot h(x,t)
\end{equation}
El diferencial de volumen se transforma como $dz = h(x,t) \, d\eta$. Al sustituir esta expresión en la integral, los límites verticales cambian de $[0, h(x,t)]$ a $[0, 1]$:
\begin{equation}
    M_{\text{finos}} = \int_{0}^{L} \int_{0}^{1} \Big( h(x,t) \cdot \phi_s(x, \eta, t) \Big) \, d\eta \, dx
\end{equation}

Si definimos el cambio de variable $Q(x, \eta, t) = h(x,t) \cdot \phi_s(x, \eta, t)$, la masa total del sistema se reduce a:
\begin{equation}
    M_{\text{finos}} = \int_{0}^{L} \int_{0}^{1} Q(x, \eta, t) \, d\eta \, dx
\end{equation}

En la formulación discreta (computacional), las dimensiones de la celda de simulación $dx$ y $d\eta$ son constantes y no varían en el tiempo ni en el espacio. Por tanto, la masa de finos contenida en una celda $(i, j)$ está dada directamente por $Q_{i,j} \, dx \, d\eta$. Al resolver la ecuación gobernante directamente para $Q$ (en lugar de para $\phi_s$), se asegura una conservación estricta de la masa a nivel de máquina, evitando que se pierda o cree sedimento artificialmente debido al estiramiento o compresión geométrica de la duna.

\section{Derivación Matemática de la Ecuación Gobernante}
Partimos de la ecuación clásica de conservación de sedimentos finos en coordenadas cartesianas 2D $(x, z)$ para un flujo incompresible ($\nabla \cdot \vec{u} = 0$):
\begin{equation}
    \frac{\partial \phi_s}{\partial t} + \frac{\partial (u \phi_s)}{\partial x} + \frac{\partial (w \phi_s)}{\partial z} = 0
\end{equation}

Queremos mapear esta ecuación al espacio transformado $(x', \eta, t')$ definido por:
\begin{align}
    x' &= x \\
    \eta &= \frac{z}{h(x,t)} \\
    t' &= t
\end{align}

Utilizando la regla de la cadena, las derivadas parciales cartesianas se expresan en función de las coordenadas transformadas como:
\begin{align}
    \frac{\partial \phi_s}{\partial z} &= \frac{\partial \phi_s}{\partial \eta} \frac{\partial \eta}{\partial z} = \frac{1}{h} \frac{\partial \phi_s}{\partial \eta} \\
    \frac{\partial \phi_s}{\partial x} &= \frac{\partial \phi_s}{\partial x'} + \frac{\partial \phi_s}{\partial \eta} \frac{\partial \eta}{\partial x} = \frac{\partial \phi_s}{\partial x'} - \frac{\eta}{h} \frac{\partial h}{\partial x} \frac{\partial \phi_s}{\partial \eta} \\
    \frac{\partial \phi_s}{\partial t} &= \frac{\partial \phi_s}{\partial t'} + \frac{\partial \phi_s}{\partial \eta} \frac{\partial \eta}{\partial t} = \frac{\partial \phi_s}{\partial t'} - \frac{\eta}{h} \frac{\partial h}{\partial t} \frac{\partial \phi_s}{\partial \eta}
\end{align}

Sustituyendo estas derivadas en la ecuación de transporte (6):
\begin{equation}
    \left( \frac{\partial \phi_s}{\partial t'} - \frac{\eta}{h}\frac{\partial h}{\partial t}\frac{\partial \phi_s}{\partial \eta} \right) + \left( \frac{\partial (u \phi_s)}{\partial x'} - \frac{\eta}{h}\frac{\partial h}{\partial x}\frac{\partial (u \phi_s)}{\partial \eta} \right) + \frac{1}{h}\frac{\partial (w \phi_s)}{\partial \eta} = 0
\end{equation}

Para eliminar los denominadores y llevar la ecuación a su forma conservativa, multiplicamos la ecuación completa por $h$:
\begin{equation}
    h \frac{\partial \phi_s}{\partial t'} - \eta \frac{\partial h}{\partial t} \frac{\partial \phi_s}{\partial \eta} + h \frac{\partial (u \phi_s)}{\partial x'} - \eta \frac{\partial h}{\partial x} \frac{\partial (u \phi_s)}{\partial \eta} + \frac{\partial (w \phi_s)}{\partial \eta} = 0
\end{equation}

Aplicando la regla del producto para derivadas, reescribimos los primeros términos de tiempo y espacio horizontal como:
\begin{align}
    h \frac{\partial \phi_s}{\partial t'} &= \frac{\partial (h\phi_s)}{\partial t'} - \phi_s \frac{\partial h}{\partial t} \\
    h \frac{\partial (u\phi_s)}{\partial x'} &= \frac{\partial (h u \phi_s)}{\partial x'} - u\phi_s \frac{\partial h}{\partial x}
\end{align}

Sustituyendo las identidades (13) y (14) en la ecuación (12) y agrupando:
\begin{equation}
    \frac{\partial (h\phi_s)}{\partial t'} + \frac{\partial (h u \phi_s)}{\partial x'} + \frac{\partial (w \phi_s)}{\partial \eta} - \left( \phi_s \frac{\partial h}{\partial t} + \eta \frac{\partial h}{\partial t} \frac{\partial \phi_s}{\partial \eta} \right) - \left( u \phi_s \frac{\partial h}{\partial x} + \eta \frac{\partial h}{\partial x} \frac{\partial (u \phi_s)}{\partial \eta} \right) = 0
\end{equation}

Dado que la altura de la duna $h$ y sus derivadas no dependen de la coordenada vertical de la simulación $\eta$, podemos introducirlas dentro de las derivadas parciales de $\eta$:
\begin{align}
    \phi_s \frac{\partial h}{\partial t} + \eta \frac{\partial h}{\partial t} \frac{\partial \phi_s}{\partial \eta} &= \frac{\partial}{\partial \eta} \left( \eta \phi_s \frac{\partial h}{\partial t} \right) \\
    u \phi_s \frac{\partial h}{\partial x} + \eta \frac{\partial h}{\partial x} \frac{\partial (u \phi_s)}{\partial \eta} &= \frac{\partial}{\partial \eta} \left( \eta u \phi_s \frac{\partial h}{\partial x} \right)
\end{align}

Reemplazando estos dos grupos en la ecuación (15) y recolectando todas las derivadas respecto a $\eta$ bajo un solo operador:
\begin{equation}
    \frac{\partial (h\phi_s)}{\partial t'} + \frac{\partial (hu\phi_s)}{\partial x'} + \frac{\partial}{\partial \eta} \left[ w\phi_s - \eta \phi_s \frac{\partial h}{\partial t} - \eta u \phi_s \frac{\partial h}{\partial x} \right] = 0
\end{equation}

Finalmente, introducimos la variable $Q = h\phi_s$:
\begin{equation}
    \frac{\partial Q}{\partial t'} + \frac{\partial (u Q)}{\partial x'} + \frac{\partial}{\partial \eta} \left[ \left( w - \eta \frac{\partial h}{\partial t} - \eta u \frac{\partial h}{\partial x} \right) \phi_s \right] = 0
\end{equation}

Esta es la forma conservativa exacta implementada en el resolvedor de diferencias finitas en diferencias de tipo upwind. El término que acompaña a $\phi_s$ dentro de la derivada vertical representa la velocidad advectiva vertical efectiva en la malla normalizada, $h w_\eta$.

\end{document}
